\section{Alternative Multi-Dimensional Academic Outcomes}
\label{sec:adv_alternative_outcomes}

While degree completion ($Y \in \{0, 1\}$) represents the primary terminal metric for institutional accreditation, financial aid affects student trajectories across continuous, multi-dimensional academic dimensions (Table \ref{tab:alternative_outcomes}).

\begin{table}[htbp]
    \centering
    \small
    \caption{Causal Treatment Effects Across Multi-Dimensional Academic Outcomes.}
    \label{tab:alternative_outcomes}
    \resizebox{\textwidth}{!}{%
    \begin{tabular}{lccccc}
        \toprule
        \textbf{Outcome Dimension} & \textbf{Cohort} & \textbf{Causal Estimand} & \textbf{Estimated ATE} & \textbf{95\% Conf. Interval} & \textbf{$p$-value} \\
        \midrule
        Continuous Sem 2 GPA & $N = 3,630$ & $\mathbb{E}[Y_1] - \mathbb{E}[Y_0]$ & \textbf{+1.478 pts} & $[+1.093, +1.863]$ & $6.69 \times 10^{-14}$ \\
        Credit Completion Ratio & $N = 3,630$ & $\mathbb{E}[Y_1] - \mathbb{E}[Y_0]$ & \textbf{+14.50\%} & $[+11.80\%, +17.20\%]$ & $1.58 \times 10^{-25}$ \\
        Full Retention (AIPW) & $N = 4,424$ & Doubly Robust AIPW & \textbf{+16.87\%} & $[+13.79\%, +19.73\%]$ & $< 1.0 \times 10^{-10}$ \\
        \bottomrule
    \end{tabular}%
    }
\end{table}

\subsection{Continuous GPA, Credit Velocity, and Cohort Retention}

On the Portuguese grading scale ($[0, 20]$), scholarship provision elevates second-semester GPA by:
\begin{equation}
\begin{aligned}
    \widehat{ATE}_{\text{GPA}} &= \mathbb{E}[Y_{\text{GPA}} \mid do(T=1)] - \mathbb{E}[Y_{\text{GPA}} \mid do(T=0)] \\
    &= \mathbf{+1.478 \text{ pts}} \quad (95\% \text{ CI}: [1.093, 1.863], \, p = 6.69 \times 10^{-14}).
\end{aligned}
\label{eq:ate_gpa}
\end{equation}
Relative to the control mean ($\mu_0 = 11.24$ pts), this represents a $+7.39\%$ surge. Furthermore, the Credit Completion Ratio (approved credits divided by enrolled credits) rises by $\mathbf{+14.50\%}$ ($95\%$ CI: $[11.80\%, 17.20\%]$, $p = 1.58 \times 10^{-25}$), directly arresting debt-driven attrition. Across the complete institutional cohort ($N = 4,424$), doubly robust AIPW yields a retention surge of $\hat{\tau}_{\text{AIPW}} = \mathbf{+16.87\%}$ ($p < 10^{-10}$), confirming that aid preserves $\approx 17$ additional enrolled students per 100 recipients.
